This section develops a structural framework for recovering the economic preferences of transformer-based language models from their generative outputs, and derives the two estimators that our identification strategy requires. A transformer model is designed to generate output autoregressively, appending one token at a time, where a token is a discrete unit of text drawn from a finite vocabulary set $\mathcal{V}$ that spans words, subwords, and characters. At each step, the model maps the accumulated input context through a learned sequence of attention and feed-forward layers to a final linear projection over $\mathcal{V}$, producing a vector of logits over all tokens in the vocabulary. Each output token is drawn stochastically from a probability distribution over $\mathcal{V}$ derived from these logits. We show for single token revealed preference problems, the logit assigned to a candidate choice acts as a cardinal index of the model's relative ranking over its revealed preferences. We formalize this mapping, define three revealed-preference indices that measure conformance to standard utility axioms on the model's forced-choice outputs, show that it admits a random utility representation nested within the multinomial logit framework of \citet{mcfadden1972conditional}, and derive the two estimators that our identification strategy requires.

The primitive of the model is a one-step forced-choice problem in which a single output token from the model constitutes a revealed preference. This structure maps naturally onto the System~1 mode of cognition described by \citet{kahneman2003maps}, where the model responds immediately and automatically from accumulated context, without deliberation, and the output is a single discrete act whose distribution over options is fully determined by the pre-softmax logits. The analogy is precise in the sense that System~1 responses are characterized as fast, associative, and context-driven, which corresponds directly to the autoregressive forward pass computed by the model with no iterative correction.\footnote{Extending the analogy to \citet{kahneman2003maps}'s System~2, one could also consider the case where the model is permitted to generate an extended chain of tokens before committing to a label. In this case the intermediate tokens function as scratchpad reasoning, allowing the model to decompose complex tradeoffs, recover dominated options through multi-step deliberation, and produce responses that can differ systematically from those elicited under single-token constraints. Recovering structural preferences from System~2 outputs raises substantively different identification challenges, as the final label token reflects not only primitive preferences but also the model's internally generated deliberation. We leave the structural analysis of System~2 choice behavior to future work and restrict attention throughout to the single-token revealed-preference setting.}

Let $x$ denote the prompt context and let $C(x) \subseteq \mathcal{V}$ denote the set of label tokens encoding the feasible options. Let $u_i(x)$ denote the raw logit on token $i \in \mathcal{V}$. Under standard softmax decoding, the model induces a next-token distribution over the full vocabulary,\footnote{For a given menu $C(x)$, the label probabilities are $\{P(i \mid x, \tau)\}_{i \in C(x)}$. Let $\bar{C}(x) = \mathcal{V} \setminus C(x)$ denote the complement; we let the residual probability $P(\varnothing \mid x, C, \tau) = \sum_{k \in \bar{C}(x)} P(k \mid x, \tau)$ collect all remaining vocabulary mass so that $\sum_{i \in C(x)} P(i \mid x, \tau) + P(\varnothing \mid x, C, \tau) = 1$.}
\begin{equation}\label{eq:softmax}
  P(i \mid x, \tau)
  = \frac{\exp(u_i(x) / \tau)}{\sum_{j \in \mathcal{V}} \exp(u_j(x) / \tau)},
  \qquad i \in \mathcal{V}.
\end{equation}
where the parameter $\tau > 0$ is temperature, which scales the dispersion of the induced choice distribution. As $\tau \downarrow 0$ the distribution concentrates all mass on the single token with the highest logit across $\mathcal{V}$. As $\tau \rightarrow \infty$ the distribution approaches the uniform over $\mathcal{V}$, spreading mass evenly across all vocabulary tokens.

\begin{definition}[Completeness Index]\label{def:completeness}
  For a prompt context $x$ and menu $C(x) \subseteq \mathcal{V}$, the \emph{completeness index} is
  \begin{equation}\label{eq:completeness}
    \mathcal{C}(x, \tau) \;\equiv\; \sum_{i \in C(x)} P(i \mid x, \tau) \;=\; \frac{\sum_{i \in C(x)} \exp\!\bigl(u_i(x)/\tau\bigr)}{\sum_{j \in \mathcal{V}} \exp\!\bigl(u_j(x)/\tau\bigr)}.
  \end{equation}
  If $C(x) \subsetneq \mathcal{V}$ and the logits on $\mathcal{V} \setminus C(x)$ are finite, then $\mathcal{C}(x, \tau) < 1$. The model satisfies \emph{exact completeness} on menu $C(x)$ at temperature $\tau$ if $\mathcal{C}(x, \tau) = 1$. When $\mathcal{C}(x,\tau) < 1$, the residual mass $1 - \mathcal{C}(x,\tau)$ falls on tokens outside the menu. Values of $\mathcal{C}$ close to one indicate that the model treats the label tokens as the relevant continuation, validating the revealed-preference interpretation of the single-token choice.
\end{definition}

We now turn to our economic interpretation of the model. The softmax rule in Equation~\eqref{eq:softmax} admits an exact random utility representation,
\begin{equation}\label{eq:rum}
  U_i = u_i(x) + \tau \epsilon_i, \qquad i \in \mathcal{V}.
\end{equation}
If we assume that the shocks $\epsilon_i$ are i.i.d.\ Type I extreme value across all $i \in \mathcal{V}$, then the softmax rule in Equation~\eqref{eq:softmax} is the full-vocabulary multinomial logit representation of the random utility model from \citep{mcfadden1972conditional}. Under this interpretation, a model's raw logit values serve as a systematic utility index, and the temperature parameter $\tau$ is the scale of the stochastic component. Consequently, we find that as $\tau \downarrow 0$, choice converges to the label with the highest logit, and the distribution recovers only the ordinal ranking of the systematic utilities (i.e. the noise vanishes and the cardinal magnitudes of utility differences become irrelevant). At $\tau = 1$, the logits enter the choice probabilities without additional rescaling, so the observed choice distribution is generated by the model's native systematic utility index. When $\tau > 1$, temperature compresses differences in that index, and as $\tau \rightarrow \infty$ the induced distribution approaches the uniform distribution over $\mathcal{V}$ that ultimately destroys the model's economic coherence.

To connect logits to structural preferences, we specify the systematic utility index as a positive affine transformation of a cardinal utility index $V_i$ over the full vocabulary,\footnote{Under expected utility, any two cardinal representations of the same preferences differ by a positive affine transformation \citep{von1947theory}. Equation~\eqref{eq:logit-spec} therefore preserves the underlying vNM preference ordering and the associated risk attitudes.}
\begin{equation}\label{eq:logit-spec}
  u_i(x) = \kappa V_i + c, \qquad i \in \mathcal{V}.
\end{equation}
The scale parameter $\kappa > 0$ maps utility differences into logit units, while the term $c$ is a common additive shifter. Substituting Equation~\eqref{eq:logit-spec} into the full-vocabulary softmax in Equation~\eqref{eq:softmax} gives
\begin{equation}\label{eq:choice-prob}
  P(i \mid x, \tau) = \frac{\exp(\frac{\kappa}{\tau} V_i)}{\sum_{j \in \mathcal{V}} \exp(\frac{\kappa}{\tau} V_j)},\qquad i \in \mathcal{V},
\end{equation}
which shows that sampled choices identify the reduced-form scale $\kappa / \tau$.

Two additional revealed-preference indices accompany the completeness index of Definition~\ref{def:completeness}. The first definition measures whether the model's logit ranking respects an exogenous dominance ordering over alternatives in the menu. The second definition, following the revealed-preference goodness-of-fit tradition of \citet{houtman1985determining} and \citet{demuynck2023computing}, measures whether the model's rankings across a collection of menus are jointly consistent with a single preference ordering over all alternatives.

% \begin{definition}[Monotonicity Index]\label{def:monotonicity}
%   For a prompt context $x$ with menu $C(x) \subseteq \mathcal{V}$ and an exogenous strict partial order $\succ_D$ on $C(x)$, let $D(x) = \{(i,j) \in C(x) \times C(x) \mid i \succ_D j\}$ denote the set of dominance pairs. Suppose $|D(x)| > 0$. The \emph{monotonicity index} is
%   \begin{equation}\label{eq:monotonicity}
%     \mathcal{M}_{adj}(x, \tau) \;\equiv\; \frac{\sum_{(i,j) \in D(x)} \frac{1}{1 + \exp\left(-\frac{u_i(x) - u_j(x)}{\tau}\right)}}{\sum_{(i,j) \in D(x)} \frac{1}{1 + \exp\left(-\frac{|u_i(x) - u_j(x)|}{\tau}\right)}}.
%   \end{equation}
%   The denominator establishes a menu-specific upper bound by evaluating the model's choices against its own absolute resolving power, $|u_i(x) - u_j(x)|$, isolating the directional adherence to economic coherence from the magnitude of the logit differences. When the model's systematic utility strictly respects the dominance ordering such that $u_i(x) > u_j(x)$ for all $(i,j) \in D(x)$, the numerator and denominator equate, yielding $\mathcal{M}_{adj}(x, \tau) = 1$. Values below one indicate that the model's logit ordering violates the maintained dominance ordering on some pairs, with the severity of the penalty scaled proportionally to the model's native confidence in that violation.
% \end{definition}

\begin{definition}[Monotonicity Index]\label{def:monotonicity}
  For a prompt context $x$ with menu $C(x) \subseteq \mathcal{V}$ and an exogenous strict partial order $\succ_D$ on $C(x)$, let $D(x) = \{(i,j) \in C(x) \times C(x) \mid i \succ_D j\}$ denote the set of dominance pairs. Suppose $|D(x)| > 0$. The \emph{monotonicity index} is
  \begin{equation}\label{eq:monotonicity}
    \mathcal{M}(x, \tau) \;\equiv\; \frac{1}{|D(x)|} \sum_{(i,j) \in D(x)} \frac{P(i \mid x, \tau)}{P(i \mid x, \tau) + P(j \mid x, \tau)}.
  \end{equation}
  As $\tau \downarrow 0$, each pairwise term converges to one if $u_i(x) > u_j(x)$ and to zero if $u_i(x) < u_j(x)$, so in the low-temperature limit $\mathcal{M}(x, \tau)$ recovers the fraction of dominance pairs for which the model's logit ordering respects $\succ_D$. The model satisfies \emph{exact monotonicity} on menu $C(x)$ if $u_i(x) > u_j(x)$ for every $(i,j) \in D(x)$, in which case the low-temperature limit of $\mathcal{M}(x, \tau)$ equals one. Values of $\mathcal{M}$ below one indicate that the model's systematic utility index does not fully respect the maintained dominance ordering, and the index quantifies the extent of the departure.
\end{definition}

\begin{definition}[Transitivity Index]\label{def:transitivity}
  Consider a collection of $N$ menus $\mathcal{D} = \{(x_t, C(x_t))\}_{t=1}^{N}$, where at each menu the model's logits induce a strict ranking $\succ_t$ over $C(x_t)$ with $i \succ_t j$ whenever $u_i(x_t) > u_j(x_t)$. A subset $S \subseteq \{1, \ldots, N\}$ is \emph{jointly consistent} if there exists a strict linear order $\succ^*$ on $\bigcup_{t \in S} C(x_t)$ such that $\succ_t \,\subseteq\, \succ^*$ for every $t \in S$. The \emph{transitivity index} for the collection $\mathcal{D}$ is
  \begin{equation}\label{eq:transitivity}
    \mathcal{T}(\mathcal{D}) \;\equiv\; \max_{S \subseteq \{1, \ldots, N\}} \frac{|S|}{N} \quad \text{subject to } S \text{ being jointly consistent.}
  \end{equation}
  For $N \geq 1$, the index is a multiple of $1/N$. The model satisfies \emph{exact transitivity} over $\mathcal{D}$ if $\mathcal{T}(\mathcal{D}) = 1$, meaning the pooled within-menu rankings across all $N$ menus admit a single preference ordering over the full set of alternatives. When $\mathcal{T}(\mathcal{D}) < 1$, some subset of menus generates a preference cycle that no total ordering can reconcile, and the index reports the largest fraction of menus that can be jointly rationalized. Under Equation~\eqref{eq:logit-spec}, exact transitivity is equivalent to the condition that the cardinal utility $V_i$ is a property of the alternative alone and does not depend on the menu context $x$ in which it appears.
\end{definition}

The three indices summarize the conditions under which the random utility representation in Equation~\eqref{eq:rum} and the invariant utility specification in Equation~\eqref{eq:logit-spec} admit a revealed-preference interpretation in one-step token choice. These diagnostics accompany the structural analysis in both observation regimes considered below. When pre-temperature logits are observed, within-menu logit differences identify the utility index directly through Equation~\eqref{eq:logit-spec}. When only sampled choices are observed, the same structure yields a likelihood based on Equation~\eqref{eq:choice-prob}. Alongside the structural estimates, both regimes report the completeness, monotonicity, and transitivity indices of Definitions~\ref{def:completeness}--\ref{def:transitivity} as revealed-preference validity diagnostics. We next apply this framework to the portfolio choice problem and derive the two structural estimators.
